% Standalone prologue; not yet integrated into the thesis.
% Build from the repository root: python3 thesis/code/build.py prologue
\documentclass[11pt]{article}
\usepackage[letterpaper,margin=0.83in]{geometry}
\usepackage{amsmath,amssymb,mathtools}
\usepackage{xcolor,graphicx}
% Collect and discard draft passages while preserving them in the source.
\NewDocumentEnvironment{comment}{+b}{}{}
\usepackage[colorlinks=true,linkcolor=blue!45!black,urlcolor=blue!45!black]{hyperref}
\definecolor{ink}{HTML}{183C50}
\definecolor{pale}{HTML}{EDF3F5}
\setlength{\parindent}{0pt}
\setlength{\parskip}{5pt}
\setlength{\headheight}{15pt}
\setlength{\emergencystretch}{2em}
\makeatletter
\def\ps@opening{%
  \def\@oddhead{\small\textcolor{ink}{TUNEM | Before the Machine}\hfil Opening for review}%
  \let\@evenhead\@oddhead
\def\@oddfoot{\footnotesize 16 September 2026 | Not yet integrated\hfil\thepage}%
  \let\@evenfoot\@oddfoot}
\makeatother
\pagestyle{opening}
\newcommand{\heading}[2]{\section*{\textcolor{ink}{#1. #2}}}
\newcommand{\refn}[1]{\textsuperscript{\hyperlink{src#1}{[#1]}}}
\newcommand{\panel}[2]{\begin{center}\setlength{\fboxsep}{8pt}\fcolorbox{ink!35}{pale}{\parbox{0.93\linewidth}{\textbf{#1}\quad #2}}\end{center}}
\newcommand{\openingfigure}[3]{%
  \begin{figure}[!htb]\centering
  \includegraphics[width=\linewidth]{figures/#1.pdf}
  \caption{\small #2}\label{#3}
  \end{figure}}
\renewcommand{\topfraction}{0.9}
\renewcommand{\bottomfraction}{0.85}
\renewcommand{\textfraction}{0.10}
\setlength{\textfloatsep}{12pt plus 2pt minus 2pt}
\newcommand{\source}[3]{\par\medskip\hypertarget{src#1}{}\textbf{[#1] #2}\par #3\par}
\newcommand{\ee}{\mathrm e}
\newcommand{\dd}{\mathrm d}
\begin{document}

{\LARGE\bfseries\color{ink} Before the Machine}\par
{\large From quarks to the physical requirements of fusion}\par
\vspace{4pt}
\textit{We will build the machine from its purpose. First we must understand the
matter we intend to change, the energy available from that change, and the
obstacle between possibility and a useful reaction rate.}

\heading{0.1}{Inside a nucleon}

A quark is one of the elementary constituents in our present theory of matter.
To locate it, begin with an atom: a positively charged nucleus surrounded by
negatively charged "quantum clouds" of electrons. Electron are also elementary constituents of matter and carry charge $-e$. Chemical reactions change electron arrangements while leaving the nuclei intact. Nuclear reactions such as fission and fusion change the nuclear arrangement itself.

Inside the nucleus are \emph{nucleons}: Protons and neutrons are baryons-- composite particles made of three quarks. Quarks come in different flavors. Two types of quarks are up quarks and down quarks. A proton can be written as $uud$ and a neutron can be written as $udd$, where
$u$ and $d$ denote up and down quarks. Quarks possess both electric charge and color charge.

The electromagnetic force is described by rotations of a circle, called the $U(1)$ group, which transforms via a matrix representation of the groups linear map $\chi(\theta) = \exp(iq\theta)$ that satisfies a circles compactification: $\chi(\theta + 2\pi) = \chi(\theta)$, which quantizes charge into an integer multiple of the smallest charge, $q \in \mathbb{Z}$. Physics doesn't tell us exactly why the electric charge is $+2e/3$ for an up quark and $-e/3$ for a down quark, anomaly cancellations give constraints, experiments fix the number: $e/3$ is the smallest positive charge we know. Given this fact, because constituent phase factors multiply (ie $e^{iq_1\theta}e^{iq_2\theta}=e^{i(q_1+q_2)\theta}$), protons have $+e$ electric charge, and
neutrons have \emph{zero net electric charge}:
\[
q_p=\left(\frac23+\frac23-\frac13\right)e=e,
\qquad
q_n=\left(\frac23-\frac13-\frac13\right)e=0.
\]

These are \emph{quark labels}: they only specify net quark content. A nucleon is a quantum state that also contains gluon fields and quark--antiquark fluctuations.\refn{1}

What binds that state? Quarks second kind of charge, called
\emph{color}. The theory of color interactions is \emph{quantum chromodynamics} (QCD), and it describes the strong nuclear force-- one of the four fundamental forces of nature. Gluons mediate the strong force and possess color charge, leading to self-interactions amongst gluons. This is unlike photons, neutrally charged electromagnetic force carriers that have no direct electric-charge self-interaction, unlike gluons with color charge.

Color is not a single number $q$ to be summed into a different state of color charge — it is a basis vector of $\mathbb{C}^3$, such as $(1,0,0)$, to be rotated via matrices whose inverse equal it's conjugate transpose and whose determinant equals to $1$, eg elements of $SU(3)$. A state is \emph{color neutral} when it is left invariant by every such rotation: for baryons this is achieved by carrying exactly one quark of each color, combined in the unique totally antisymmetric combination, $\epsilon_{ijk}q^iq^jq^k$, that $SU(3)$ leaves unchanged. This explains why three quarks can form a color-neutral state, whereas two or four quarks alone cannot: three colors demand three quarks to form this invariant, fully analogous to how a determinant vanishes unless you antisymmetrize over as many vectors as there are dimensions. Indeed $\epsilon_{ijk}U^i{}_aU^j{}_bU^k{}_c=(\det U)\,\epsilon_{abc}$, so if color rotations were allowed any determinant other than $1$, this three-quark combination would change by that factor and no longer be neutral; that is why the group is $SU(3)$.\refn{2}

Most nucleon mass is not the sum of the small up- and down-quark mass parameters.
It comes from the energy of the interacting quark--gluon system. Lattice QCD,
which approximates the fields on a spacetime grid and controls extrapolations,
has reproduced the light-hadron mass spectrum. Most ordinary matter's mass is
in nucleons, predomenantkly made of dynamical QCD energy.\refn{3}

\openingfigure{figN0_atom_to_quarks}{Inside a helium atom is the nucleus, inside the nucleus are the nucleons: protons and neutrons. Inside a nucleon such as a proton are three quarks.  Colored regions are schematic; the proton's
$uud$ reflect two up-flavored quarks and one down quark. Neutrons follow $udd$.}{fig:n0}

\heading{0.2}{Confinement, and the boundary of the explanation}

QCD has an unusual scale dependence. Quarks are never observed isolated in nature. Nucleons are three quarks grouped together, an example of baryons. More generally, any number of quarks grouped together are called hadrons.

When quarks are extremely close together, the strong interaction between them becomes weak, and they behave almost like free particles. This is called \emph{asymptotic freedom}. However, as quarks are pulled farther apart, their color field becomes concentrated into a narrow tube of gluon field called a \emph{flux tube}. Stretching this tube requires more and more energy. Eventually, enough energy is stored in the tube that it becomes energetically favorable to create a new quark--antiquark pair precisely when the flux tube breaks.  In conclusion, no isolated quark is produced: each separated quark instead joins with one of the newly created partners to form a new color-neutral particle. In this way, quarks are never observed in isolation. This is called \emph{quark confinement} or \emph{color confinement}.\refn{4}

The flux-tube picture is supported by QCD calculations and lattice simulations, but mathematically proving color confinement analytically from four-dimensional QCD remains an open problem. The key feature is constant linear growth: once the flux tube has formed, each additional unit of separation stores approximately the same additional amount of field energy. This linear confinement seems to produce an
\emph{area law}. A simple model keeps the tube's thickness, or cross section, roughly fixed as its length changes. The quark separation (r) therefore parametrizes the constant linear growth of the potential energy.

\[
V(r)\simeq V_0+\mathcal T_s r,
\qquad
F_r=-\frac{\dd V}{\dd r}\simeq-\mathcal T_s.
\]

The first step in proving such a physical gap is proving a simpler mathematical formalism. Strip away the quarks and keep only the gluon field, so no tube can ever break.
Light shows what a force field can do without a gap: a long enough
electromagnetic wave carries arbitrarily little energy, $E=hc/\lambda$. Does the
self-interacting gluon field allow the same, energies arbitrarily close to zero?
Or is there a smallest positive energy $\Delta>0$, so that its lightest
excitations are massive bundles of field rather than long waves? This possible empty interval above the vacuum is the \emph{Yang--Mills mass gap}, one of the deepest open problems in mathematical physics.

\panel{A Millennium Problem}{The Yang--Mills
existence and mass-gap problem asks for a rigorous construction of the
four-dimensional quantum theory and a strictly positive energy gap above its
vacuum. Writing $H$ for its energy operator and setting vacuum energy to zero,
the gap condition is
\[
\operatorname{spec}(H)\cap(0,\Delta)=\varnothing
\quad\hbox{for some }\Delta>0.
\]}

Simulations suggest that no physical excitation can have an energy strictly between zero and \(\Delta\). Providing a conceptual mathematical proof is expected to shed deep insight into the structure of the universe. Confinement is the physical mechanism that parametrizes energy by the length of flux tubes; the mass gap may be understood as a question about the spectrum \(\operatorname{spec}(H)\) of those energy levels. Mathematically, the spectrum is the set of all \(\lambda\) for which the operator \(T-\lambda I\) does not have a bounded inverse. This generalizes the concept of eigenvalues from finite-dimensional vector spaces to the infinite-dimensional \emph{Hilbert spaces} found in quantum mechanics.\refn{6}

Gluons can themselves form bound states called \emph{glueballs}.
In 2026, the BESIII experiment in China experimentally confirmed the existance of glueballs via the $X(2370)$ particle, giving experimental
substance to the idea that the carriers of a force can organize into
matter.\refn{20}

We now have the constituents of a nucleus. The next question is what happens
when nucleons meet. How can two color-neutral particles attract, and what
determines whether their shared arrangement releases energy? We will follow
that question from the interaction itself to a fusion reaction.

% Author's speculative story: retained unchanged for later revision.
\begin{comment}
\textbf{A speculative picture.} The \emph{liquid-drop model} offers a potential bridge from confinement to nuclear binding to fusion. It pictures the nucleus
as a self-bound liquid drop. The drop can deform, vibrate, and respond collectively, but deforming its surface costs energy.\refn{19}
Could an analogous language of collective stiffness help us picture how the
quark--gluon field confines its constituents and develops a spectrum of massive
excitations?

Imagine two nucleons close together. Their internal fields and fluctuations
participate in a shared interaction 
described through meson exchange and other effective couplings inherited from
QCD.\refn{8} What reaches beyond each nucleon is its capacity to correlate with
its neighbours. When the combined state has lower energy than the separated
nucleons, the difference is nuclear binding energy. Within the resulting drop,
surface vibrations can themselves mediate additional attraction, as developed
in \emph{The Nuclear Cooper Pair}.\refn{19} This suggests a succession of
organization: quark--gluon dynamics builds nucleons, interactions among nucleons
build nuclei, and collective motion reshapes their binding. Fusion opens a path
between nuclear arrangements; when the products are more tightly bound, the
energy difference emerges as motion or radiation. For deuterium and tritium,
that difference is about $17.6\,\mathrm{MeV}$.\refn{11} The question inviting
further thought is how this sequence produces such different scales: the large
rest energy of a nucleon, the smaller binding energy of a nucleus, and the
collision energies at which fusion becomes appreciable.

Ordinary liquids suggest a provocative electromagnetic comparison: their
cohesion emerges from the charge structure and quantum correlations of their
constituents. Rutherford's 1920 conjecture pictured a neutron as a proton
tightly combined with an electron; a distinct modern pion-cloud description
includes a virtual proton accompanied by a negatively charged pion.\refn{21}\refn{22}
Both images prompt a question about how internal charge structure can allow a
neutral composite to interact with its surroundings. How far could polarization
and quantum exchange take an electromagnetic account of nuclear cohesion, and
could collective response help us understand the barrier through which fusing
nuclei approach? The \emph{Kondo effect} supplies a further way to think about
how electromagnetic phenomenon can mirror nuclear phenomenon such as asymptotic freedom: an antiferromagnetic exchange interaction between a magnetic
impurity and a metal's electrons grows strong as the energy scale decreases,
producing collective screening below a characteristic Kondo scale.\refn{7}
This also invites us to ask whether similar patterns of scale evolution can illuminate
the hierarchy from color confinement to nuclear binding, and how the surrounding
medium influences access to a fusion reaction. In the next section we focus on more standard perspectives.
\end{comment}

\openingfigure{figN1_string_breaking}{Static heavy sources $Q,\bar Q$ and
creation of a light $q,\bar q$ pair. The final groups are color neutral; dashed
ovals group constituents rather than marking physical walls. The qualitative
potential sketch omits the smooth mixing near string breaking and has no
physical scale. Original schematic based on source 5; not a proton model.}{fig:n1}

\clearpage
\heading{0.3}{How can two nucleons attract?}

Start with one proton and one neutron, far apart. Bring them closer.
At separations comparable to their own sizes, their quark--gluon structures
participate in an interaction that can hold the pair together. The resulting
bound pair is called a \emph{deuteron}: the nucleus of deuterium, one of our
fusion fuels. Its existence gives us a concrete question to explain:
what makes the pair energetically preferable to its separated parts?

\textbf{First meet the pion.}
QCD can organize quarks into several kinds of color-neutral particle.
We have met the three-quark content of a proton or neutron. A quark can also
combine with an \emph{antiquark}. An antiquark is the corresponding antimatter
constituent: it has the same mass and opposite electric charge, and carries
the conjugate kind of color charge. A quark--antiquark combination can be
color neutral. Such combinations form the ordinary \emph{mesons}.

The lightest mesons are \emph{pions}. Their dominant quark content involves
the up and down flavors we already know. For example, an up quark and a
down antiquark give the positively charged pion:
\[
\pi^+:\quad u\bar d,\qquad
q_{\pi^+}=\frac{2e}{3}+\frac{e}{3}=e.
\]
The bar denotes an antiquark. Reversing the flavors gives $\pi^-$ with
charge $-e$; the neutral pion, $\pi^0$, involves a quantum combination of
$u\bar u$ and $d\bar d$. Gluon fields and further quantum fluctuations are
part of each pion's full state, just as they are part of a nucleon's.\refn{26}

Where do these particles come from? The quark and gluon fields can exchange
energy and create quark--antiquark pairs. In a collision with enough available
energy, those fields can organize into a pion that emerges as an observable
particle. A pion is therefore another form of matter that the same QCD
dynamics can build. Its unusually small mass has a deeper symmetry explanation,
developed in the optional note at the end of this opening.

\textbf{How does a pion connect two nucleons?}
An interaction between nucleons transfers momentum and changes their joint
energy. One important contribution can be calculated as one nucleon emitting
a pion and the other absorbing it. Here the pion is an \emph{intermediate}
part of the quantum process. This is the meaning of \emph{virtual pion
exchange}: the calculation includes the pion internally, while the incoming
and outgoing particles can both be nucleons.

This description follows from QCD by changing the level of detail. At the
energies relevant to nuclear binding, we use nucleons and pions as the working
objects and let their interactions summarize the underlying quark--gluon
dynamics. An \emph{effective field theory} makes that procedure systematic:
QCD's symmetries constrain the allowed interactions, measurements determine
their strengths, and successive terms improve the calculation. Single-pion
exchange supplies the longest-reaching part of this strong interaction;
closer in, further contributions become important.\refn{8}

\textbf{Why does this interaction have such a short reach?}
We need two units to see the scale. A femtometre is
$1\,\mathrm{fm}=10^{-15}\,\mathrm m$. An electronvolt is the energy acquired
by a charge of magnitude $e$ crossing one volt:
$1\,\mathrm{eV}=1.602176634\times10^{-19}\,\mathrm J$.
A kiloelectronvolt is $10^3\,\mathrm{eV}$ and a megaelectronvolt is
$10^6\,\mathrm{eV}$. We write them as keV and MeV.\refn{12}

The pion's mass is the key. In the simplest massive-exchange description,
the distance dependence contains
\[
\frac{e^{-r/\ell_\pi}}{r},
\qquad
\ell_\pi=\frac{\hbar}{m_\pi c}
        =\frac{\hbar c}{m_\pi c^2}.
\]
Here $r$ is separation, $m_\pi$ is the pion mass, $c$ is the speed of light,
and $\hbar=h/(2\pi)$ is the reduced Planck constant. This is the
\emph{Yukawa range}: the massive-field equation produces exponential decay.
The full pion-exchange interaction also depends on the nucleons' spin and
on whether they are protons or neutrons.\refn{24}

The measured pion rest energy $m_\pi c^2$ is about $135$--$140\,\mathrm{MeV}$.
Using $\hbar c\simeq197.3\,\mathrm{MeV\,fm}$, the range is therefore
\[
\ell_\pi\simeq
\frac{197.3\,\mathrm{MeV\,fm}}{140\,\mathrm{MeV}}
\simeq1.4\,\mathrm{fm}.
\]
Every additional distance $\ell_\pi$ multiplies the exponential factor by
$e^{-1}\simeq0.37$. That is a rapid fading with distance. Nuclear attraction
can be substantial when the particles are close, while electrical repulsion
between protons remains important much farther out. We will meet this
competition again when we try to bring two fuel nuclei together.

\heading{0.4}{Why does a nucleus have a finite size?}

Attraction alone leaves something unexplained. If bringing the particles
closer lowers their energy, what stops a nucleus from shrinking indefinitely?
The answer brings quantum motion into the story.

A quantum particle occupies a state spread over space. We describe that state
by a wavefunction, a function whose squared magnitude gives the probability
density for position. Concentrating that probability into a smaller region
requires a wider spread of momenta. Momentum carries kinetic energy, so
localizing the particle has an energy cost.

\textbf{A small calculation makes the competition visible.}
Look along one coordinate $x$. Write $\Delta x$ for the standard deviation of
the position distribution and $\Delta p_x$ for that of the corresponding
momentum. These measure the spreads of possible measurement outcomes in the
same quantum state. The uncertainty relation is\refn{23}
\[
\Delta x\,\Delta p_x\geq\frac{\hbar}{2}.
\]
Let $L=\Delta x$ be our measure of spatial width. Dividing by $L$ gives
$\Delta p_x\geq\hbar/(2L)$.

For a particle of mass $m$, the nonrelativistic kinetic energy associated
with this direction is $p_x^2/(2m)$. Angle brackets denote its quantum average.
The definition of variance gives
\[
\langle p_x^2\rangle=(\Delta p_x)^2+\langle p_x\rangle^2.
\]
The smallest cost occurs when the mean momentum is zero. Combining these
relations, step by step, gives
\[
\langle K_x\rangle
=\frac{\langle p_x^2\rangle}{2m}
\geq\frac{(\Delta p_x)^2}{2m}
\geq\frac{1}{2m}\left(\frac{\hbar}{2L}\right)^2
=\frac{\hbar^2}{8mL^2}.
\]
The important dependence is $1/L^2$: halving the width quadruples this
minimum kinetic-energy cost. This one-coordinate calculation isolates the
cost of localization; a nuclear calculation also includes the other
coordinates, interactions, and the motion of all its constituents.

To see its size, use a nucleon mass $m_Nc^2\simeq939\,\mathrm{MeV}$ and
choose $L=1\,\mathrm{fm}$. Multiplying numerator and denominator by $c^2$
lets us use the units already introduced:
\[
\frac{\hbar^2}{8m_NL^2}
=\frac{(\hbar c)^2}{8(m_Nc^2)L^2}
\simeq\frac{(197.3\,\mathrm{MeV\,fm})^2}
{8(939\,\mathrm{MeV})(1\,\mathrm{fm})^2}
\simeq5.2\,\mathrm{MeV}.
\]
Even this lower bound is on a nuclear energy scale. At half that width it
becomes about $21\,\mathrm{MeV}$.

\textbf{Binding is the outcome of competing contributions.}
Measure the energy of a proton--neutron pair relative to the same particles
separated and at rest. Its internal energy contains a positive kinetic
contribution and an interaction contribution that can be negative:
\[
E_{\rm relative}=\langle K_{\rm internal}\rangle+\langle V_{\rm interaction}\rangle.
\]
Here $V_{\rm interaction}$ describes their nuclear interaction, and the
averages are taken in the same two-particle state.

A very extended state samples little of the attractive region. A very
concentrated state pays a high kinetic cost; commonly used nuclear potentials
also have a strongly repulsive short-distance core. Between those extremes,
the pair may find a lower-energy state with a finite spatial spread. A
calculation with the actual interaction determines whether that state is
bound and by how much.\refn{8}\refn{24}

For the deuteron, measurements give
$E_{\rm relative}\simeq-2.225\,\mathrm{MeV}$. It takes
$2.225\,\mathrm{MeV}$ to separate it into a proton and a neutron with no
remaining relative motion. Its \emph{binding energy} is that positive
separation cost. Conversely, forming the deuteron from free particles releases
energy into other products, such as an emitted photon.\refn{11}

\textbf{What changes when we add more nucleons?}
There are more attractive interactions, but also more states to accommodate.
The Pauli exclusion principle says that two identical fermions cannot occupy
the same complete one-particle state. Fermions are particles with half-integer
spin; protons and neutrons both belong to this class. Protons obey the
exclusion rule among themselves, as do neutrons. A state includes both its
spatial wavefunction and its \emph{spin}, the particle's intrinsic angular
momentum.\refn{25}

For a nucleon, spin has two possible projections along any chosen axis.
In an independent-particle description, a spatial orbital---one available
spatial wavefunction---can therefore hold two protons with opposite spin
projections, and likewise two neutrons. As more identical nucleons are added,
further spatial orbitals must be used. Their energies and the interactions
among all the particles determine the binding of the new nucleus.

We can now see what nuclear binding asks of the physics: find an allowed
shared quantum state whose total energy lies below the separated-particle
energy. The next step is to compare two different nuclear arrangements.
That comparison will give us the energy available from fusion.

\heading{1}{Reading the reaction before using it}

Let $Z$ count protons and $N$ count neutrons; $A=Z+N$ counts nucleons. The symbol
${}^{A}_{Z}X$ specifies a nucleus. Changing $Z$ changes the element $X$; changing
$N$ at fixed $Z$ gives a different isotope. For our fuel and products,
\begin{center}
\begin{tabular}{@{}llll@{}}
\hline
Nucleus & Symbol & Constituents & Electric charge\\
\hline
Deuterium (deuteron) & $D={}^{2}_{1}\mathrm H$ & $1p+1n$ & $+e$\\
Tritium (triton) & $T={}^{3}_{1}\mathrm H$ & $1p+2n$ & $+e$\\
Helium-4 (alpha) & $\alpha={}^{4}_{2}\mathrm{He}$ & $2p+2n$ & $+2e$\\
Free neutron & $n$ & unbound neutron & $0$\\
\hline
\end{tabular}
\end{center}
The principal D--T fusion channel is
\begin{equation}
D+T\longrightarrow\alpha+n.\label{eq:reaction}
\end{equation}
Its bookkeeping is $2+3=4+1$ nucleons and $1+1=2+0$ units of charge.
The neutron on the right is not inside the helium nucleus. Energy, momentum,
and angular momentum must also balance, but these conditions do not force an
encounter to fuse: the nuclei can scatter without changing species.\refn{9}

The Sun's proton--proton route includes $p+p\to D+e^++\nu_e$, where the weak
interaction converts a proton into a neutron during formation of deuterium;
$e^+$ is a positron and $\nu_e$ an electron neutrino.
Our D--T reaction already has neutron-containing fuel and needs no such weak
conversion.\refn{10}

\heading{2.1}{The energy is a difference between arrangements}

The deuteron gave us the basic idea: binding lowers the energy of a shared
arrangement. We can measure that lowering by weighing the arrangement.
Einstein's relation $E_{\rm rest}=mc^2$ connects the mass of an object to
its energy when the object as a whole is at rest. Its internal motion and
interactions are included in this rest energy.

Imagine separating a nucleus into its $Z$ protons and $N$ neutrons. Supply
just enough energy to leave those particles far apart with no kinetic energy.
The energy supplied is the binding energy $B$. Energy conservation then reads
\[
\underbrace{m(Z,N)c^2}_{\text{nucleus at rest}}+B
=\underbrace{Zm_pc^2+Nm_nc^2}_{\text{separated nucleons at rest}}.
\]
Here $m_p$ and $m_n$ are the proton and neutron masses, and $m(Z,N)$ is the
mass of the nucleus alone; electrons are left out of both sides. Solving for
the separation cost gives
\begin{equation}
B(Z,N)=\bigl[Zm_p+Nm_n-m(Z,N)\bigr]c^2.\label{eq:binding}
\end{equation}
Thus a bound nucleus weighs less than its separated constituents by $B/c^2$.
This difference is the \emph{mass defect}. Greater binding means a larger
separation cost, and therefore a lower rest energy for that inventory of
nucleons.\refn{9}

\textbf{Now compare the two arrangements in our reaction.}
The energy made available by one D--T reaction, called its $Q$ value, is the
initial rest energy minus the final rest energy. For each bound nucleus we
substitute ``separated-nucleon energy minus binding energy.'' The free neutron
on the right contributes its own rest energy:
\begin{align}
Q&=(m_D+m_T-m_\alpha-m_n)c^2\notag\\
 &=[(m_p+m_n)c^2-B_D]+[(m_p+2m_n)c^2-B_T]\notag\\
 &\quad-[(2m_p+2m_n)c^2-B_\alpha]-m_nc^2\notag\\
 &=B_\alpha-B_D-B_T.\label{eq:q}
\end{align}
Both sides contain two protons and three neutrons, so their separate rest
energies cancel exactly. What survives is how much more tightly the final
arrangement is bound. Evaluated mass data give, rounded to
$0.001\,\mathrm{MeV}$, $B_D\simeq2.225$, $B_T\simeq8.482$, and
$B_\alpha\simeq28.296\,\mathrm{MeV}$. Thus
\[
Q\simeq28.296-2.225-8.482=17.589\,\mathrm{MeV}
\simeq17.6\,\mathrm{MeV}.\refn{11}
\]
These three binding energies are measured inputs. The calculation shows how
they determine the available fusion energy. A microscopic nuclear calculation
has a further job: finding the interacting quantum states that reproduce
these measured values.

The alpha particle makes the gain possible. In a simple orbital picture,
its two protons can occupy one spatial orbital with opposite spin projections,
and its two neutrons can do the same. The exclusion rule therefore permits a
compact four-nucleon arrangement. Its actual binding follows from the balance
of nuclear attraction, quantum motion, and electrical repulsion; the measured
$28.296\,\mathrm{MeV}$ tells us how favorable that balance is.\refn{9}\refn{25}

\textbf{Where does the released energy go?}
Into the motion of the alpha particle and neutron. To calculate the sharing,
choose a frame in which the total momentum of the incoming pair is zero.
This is the \emph{center-of-momentum frame}. Conservation gives
\[
\mathbf p_\alpha+\mathbf p_n=0,
\qquad |\mathbf p_\alpha|=|\mathbf p_n|=p.
\]
The products recoil in opposite directions. At the same momentum, the lighter
particle moves faster and carries more kinetic energy. Using the
nonrelativistic expression $K=p^2/(2m)$ makes the ratio explicit:
\[
\frac{K_n}{K_\alpha}
=\frac{p^2/(2m_n)}{p^2/(2m_\alpha)}
=\frac{m_\alpha}{m_n}\simeq4.
\]
For incoming kinetic energy small compared with $Q$, the final kinetic
energies add to approximately $Q$. Substituting $K_n\simeq4K_\alpha$ gives
\[
K_\alpha+K_n\simeq Q
\quad\Longrightarrow\quad 5K_\alpha\simeq Q.
\]
Thus $K_\alpha\simeq Q/5\simeq3.5\,\mathrm{MeV}$ and
$K_n\simeq4Q/5\simeq14.1\,\mathrm{MeV}$. This leading estimate will matter
when we build the reactor: most of the released energy travels with a neutral
particle. Relativistic two-body kinematics refines these values and includes
the incoming kinetic energy.\refn{13}

\openingfigure{figN3_binding_recoil}{The same $2p+3n$ inventory is rearranged.
The two final recoil arrows have equal momentum magnitude, not equal kinetic
energy; energies use the leading nonrelativistic estimate in the text. Rest
energies are measured relative to separated nucleons. Nucleon circles encode
counts, not exact positions or nuclear shapes. Original diagram using sources
11 and 13 and Eq.~\eqref{eq:q}.}{fig:n3}

\heading{2.2}{An energetic downhill route can still have a barrier}

The wider pattern is visible by plotting binding energy \emph{per nucleon},
$B/A$, against nucleon count $A$. It rises broadly from light nuclei to the
iron--nickel region and decreases toward very heavy nuclei. Many light-nucleus
fusion reactions and heavy-nucleus fission reactions therefore release energy:
their products have greater \emph{total} binding energy than their reactants.
For any proposed reaction, Eq.~\eqref{eq:q}'s method gives the decisive test:
compare the total initial and final rest energies.\refn{9}

Why a peak? The effective nuclear attraction is short ranged: a nucleon mostly
benefits from nearby neighbours. Small nuclei have a large surface fraction,
so many nucleons have fewer neighbours; larger nuclei initially gain more binding
per particle. But electric repulsion acts between every pair of protons, making
large charged nuclei increasingly costly. This competition gives the broad
trend. Quantum shell structure, proton--neutron imbalance, and pairing determine
the detailed curve; this argument alone does not calculate the peak.\refn{9}

The top of the curve is a broad iron--nickel neighborhood.
${}^{62}\mathrm{Ni}$ has the largest measured $B/A$, about
$8.795\,\mathrm{MeV}$; ${}^{56}\mathrm{Fe}$ is nearby at about
$8.790\,\mathrm{MeV}$. Stability against a particular decay still depends on
its possible products and conserved quantities.\refn{11}

\textbf{The energy is downhill; the entrance is uphill.}
The final arrangement is favorable, but the nucleons must first get close
enough to rearrange. A deuteron and triton both carry $+e$. Their electrical
repulsion begins far outside the region where nuclear attraction is strong.
As they approach, some of their kinetic energy becomes electrical potential
energy, slowing the approach.

Let $r$ be the distance between the two nuclei. Treating them as point
charges outside their overlap region, Coulomb's law gives the outward force
$F_r=e^2/(4\pi\epsilon_0r^2)$. The constant $\epsilon_0$ is the vacuum
permittivity, which sets the electrical-force scale in SI units. The potential
energy $U$ satisfies $F_r=-\dd U/\dd r$. Choose zero energy at infinite
separation and integrate the work needed to push inward:
\begin{equation}
U(r)=\int_r^\infty\frac{e^2}{4\pi\epsilon_0s^2}\,\dd s
=\frac{e^2}{4\pi\epsilon_0}\left[-\frac1s\right]_r^\infty
=\frac{e^2}{4\pi\epsilon_0r}.\label{eq:coulomb}
\end{equation}
Using $e^2/(4\pi\epsilon_0)\simeq1.440\,\mathrm{MeV\,fm}$, at
$r=5\,\mathrm{fm}$ the cost is
\[
U(5\,\mathrm{fm})\simeq\frac{1.440\,\mathrm{MeV\,fm}}{5\,\mathrm{fm}}
=0.288\,\mathrm{MeV}=288\,\mathrm{keV}.\refn{12}
\]
For comparison, a head-on encounter with only $10\,\mathrm{keV}$ of relative
kinetic energy reaches its classical turning point when $U(r)=10\,\mathrm{keV}$.
Solving gives $r\simeq1.440/0.010=144\,\mathrm{fm}$: much farther out than
the few-femtometre scale of nuclear attraction. The incoming pair would turn
around before reaching the region where fusion can occur.

\textbf{A quantum wave can reach farther.}
The wavefunction describes a spread of possible separations. Its behavior
inside a barrier follows from the Schr\"odinger equation. We can see the
mechanism with an ordinary differential equation by replacing the curved
barrier with a constant height $U_0$ over a finite interval of width $d$.
Write $x$ for position in this toy model, $E<U_0$ for the conserved energy,
and $\mu$ for the particle's mass. The stationary equation is\refn{14}
\[
-\frac{\hbar^2}{2\mu}\frac{\dd^2\psi}{\dd x^2}+U_0\psi=E\psi.
\]
Subtract $U_0\psi$ and multiply by $-2\mu/\hbar^2$:
\[
\frac{\dd^2\psi}{\dd x^2}
=\frac{2\mu(U_0-E)}{\hbar^2}\psi
=\chi^2\psi,
\qquad \chi=\frac{\sqrt{2\mu(U_0-E)}}{\hbar}>0.
\]
The solutions are exponentials, $\psi(x)=C\ee^{\chi x}+D\ee^{-\chi x}$.
At each edge of the barrier, matching the wavefunction and its derivative
joins this solution to the waves outside. A finite transmitted wave remains
on the far side. This is \emph{tunnelling}. For a thick barrier the
transmitted probability is suppressed approximately by the factor
$\ee^{-2\chi d}$: penetration becomes easier when the barrier is thinner
or when the incoming energy is closer to its height.

For the approaching nuclei, the corresponding coordinate is their separation,
and its kinetic term uses the reduced mass
$\mu=m_Dm_T/(m_D+m_T)$. The actual calculation includes the varying Coulomb
potential and the possible nuclear outcomes. The toy model isolates the
physical opening: a wave can penetrate to small separations even when a
classical trajectory would turn back. We still need the probability of
fusion once the nuclei interact, which is our next question.

\openingfigure{figN4_binding_and_barrier}{Different horizontal axes, different
questions. (A) AME2020 entries with measured binding values and $A\leq260$;
estimated values marked \texttt{\#} are excluded. The line is an upper envelope
of the discrete entries, not a reaction path. Arrows indicate broad possible
energy-releasing trends, not that every channel is favorable. (B) An exact
stationary solution of the stated constant-barrier toy model: $E/U_0=0.35$,
width $2.2\ell$, with $\ell=\hbar/\sqrt{2\mu U_0}$. Matching $\psi$ and its
derivative gives nonzero transmission; the lower curve is amplitude magnitude,
not an energy curve. Original calculation; sources 11 and 14.}{fig:n4}

\heading{3}{From a possible reaction to a probability}

The symbol $D+T\to\alpha+n$ identifies an outcome. To predict its frequency,
we need a quantum amplitude. A quantum state can be represented by a vector;
time evolution changes that vector according to the interactions. For
scattering, an operator $S$ connects specified incoming states to outgoing
states. In a simplified orthonormal wave-packet basis, write
\[
S|i\rangle=\sum_f a_f|f\rangle,
\qquad a_f=\langle f|S|i\rangle,
\qquad P_f=|a_f|^2.
\]
$|i\rangle$ denotes the prepared incoming state, $|f\rangle$ a resolved outgoing
state, and $a_f$ its complex coefficient. Orthogonality gives
$\langle f|g\rangle=0$ for distinct basis states and $1$ when $f=g$;
probability conservation then reads $\sum_f|a_f|^2=1$. Actual continuum momenta
require integrals or finite detector bins rather than a probability for one
exact momentum.\refn{15}

An amplitude includes both magnitude and phase. Contributions to the same
indistinguishable outcome add \emph{before} squaring; interference can therefore
increase or suppress that outcome. We do not select nuclear products by counting
nucleons alone.\refn{15}

\textbf{How would we measure this?} Send a narrow-energy deuteron beam toward a
thin target containing $N_T$ tritons. Let $\Phi_D$ be the incident number flux,
in $\mathrm{m^{-2}s^{-1}}$, and let $\dot N_{\rm fus}$ count fusion events per
second. In the dilute, independent-target limit, the operational definition is
\begin{equation}
\sigma_{DT}(E)=\frac{\dot N_{\rm fus}}{\Phi_DN_T}.
\label{eq:sigma}
\end{equation}
Its units are area. It is the reaction's \emph{effective} area at relative
collision energy $E$, not the literal cross-sectional area of a little hard
nucleus. The thin-target assumption avoids appreciable beam depletion and
energy loss. Connecting an $S$-matrix element to this area requires incoming
flux normalization, outgoing phase-space integration, and appropriate spin
sums or averages; one cannot simply set $\sigma=|a|^2$.\refn{13}

Barrier penetration is only part of the amplitude. D--T fusion is enhanced by
a resonance, a temporarily correlated five-nucleon configuration with the
quantum numbers of unbound helium-5. A microscopic calculation by Hupin and
colleagues treats incoming $D+T$ and outgoing $\alpha+n$ channels with nuclear
interactions from effective field theory. Its high-accuracy comparison includes
a small phenomenological resonance adjustment. This gives a concrete bridge
from a microscopic nuclear model to measured reaction probabilities, with the
model's assumptions still visible.\refn{16}

\panel{Where do path integrals enter?}{An amplitude can also be built by
adding contributions from possible histories. Each history carries a phase
set by its action divided by $\hbar$. This gives another way to investigate
barrier penetration: in a suitable approximation, special paths in imaginary
time reveal the tunnelling suppression. The optional instanton note develops
this connection after the main story.\refn{17}\refn{28}}

\heading{4}{From one encounter to hot fuel}

A reactor needs many reactions, not merely a possible one. Let $n_D$ and $n_T$
denote numbers of deuterons and tritons per unit volume. First imagine all
encounters have relative speed $v$. In a short time $\dd t$, one deuteron
samples an effective reaction volume $\sigma_{DT}v\,\dd t$. That volume contains,
on average, $n_T\sigma_{DT}v\,\dd t$ tritons. Multiplying by $n_D$ deuterons
per unit volume gives the event rate density
$\mathcal R=n_Dn_T\sigma_{DT}v$.

Real particles do not all move at the same speed. Let $g(v)\,\dd v$ be the
probability that a randomly selected D--T pair has relative speed in
$[v,v+\dd v]$, with $\int_0^\infty g(v)\,\dd v=1$. Averaging the preceding
argument gives
\begin{equation}
\langle\sigma v\rangle=\int_0^\infty\sigma_{DT}(v)v\,g(v)\,\dd v,
\qquad \mathcal R=n_Dn_T\langle\sigma v\rangle.\label{eq:rate}
\end{equation}
The brackets mean an average, not multiplication of separately averaged factors.
Here $\sigma_{DT}(v)$ is the same cross section, indexed by relative speed
rather than collision energy. For unlike species there is no factor $1/2$:
each D--T pair is counted once.
This is a local, dilute, binary-collision model with uncorrelated incoming
particles; it is not a formula for an arbitrary strongly correlated state.
Dimensional checking gives
$[\langle\sigma v\rangle]=\mathrm{m^3s^{-1}}$ and
$[\mathcal R]=\mathrm{m^{-3}s^{-1}}$. Multiplying by the energy per event gives
the local \emph{nuclear} power density, $\mathcal P_{\rm fus}=Q\mathcal R$,
not net electrical output.

For near-thermal fuel, ion temperature sets the velocity distribution. We keep
temperature in kelvin, $T_i$, distinct from its energy scale,
$\Theta_i=k_BT_i$, where $k_B$ is Boltzmann's constant. For example,
\[
\Theta_i=10\,\mathrm{keV}=10^4\,\mathrm{eV}
\quad\Rightarrow\quad
T_i=\frac{10^4}{8.617333262\times10^{-5}}\,\mathrm K
\simeq1.16\times10^8\,\mathrm K.\refn{12}
\]
This is a scale illustration, not a temperature optimum derived from the
Coulomb barrier. Evaluated cross sections, energy distributions, and losses
must be combined to choose an operating range. The familiar multi-keV D--T
regime will be justified by those calculations, not assumed to follow from
``hotter is better.''\refn{18}

In the hot-fuel regime used for magnetic fusion, electrons no longer remain
bound to individual nuclei: the fuel is an ionized plasma of positive nuclei
and negative electrons. It is not a quark--gluon plasma;
nucleons remain confined composite particles.\refn{1}
Positive and negative charge can balance overall even though each constituent
responds to electromagnetic fields.

\panel{The first machine question}{How can we give the fuel enough energetic
encounters, and retain that energy long enough, for fusion power to compete
with losses? In the magnetic-confinement route pursued here, sustained contact
between the hot core and material walls would rapidly remove energy and add
impurities. We need separation of core fuel from material surfaces. We have not
yet earned a torus or a coil: next we ask whether electromagnetic forces can
provide the required control.}

The charged alpha particle and neutral neutron will later demand different
energy-handling strategies. Microscopic color confinement and macroscopic
plasma confinement are different problems; the latter does not follow from
the former by analogy.

\section*{\textcolor{ink}{Optional deeper look: why are pions so light?}}

A pion has about one-seventh of a nucleon's mass. That small mass gave the
nuclear interaction its longest reach in Section~0.3. The deeper explanation
is a connection between symmetry and low-energy motion.

Picture an energy landscape whose lowest points form a circular trough.
Moving radially away from the trough costs energy; moving around its bottom
does not. If the system settles at one point on the circle, its chosen state
has less symmetry than the landscape: rotating the landscape changes nothing,
but rotating the chosen state moves it to a different point. This is the
basic idea of \emph{spontaneous symmetry breaking}.

In a field theory, neighboring regions can sit at slightly different points
along such a family of equivalent ground states. Slowly varying disturbances
then cost little energy. Under the conditions of Goldstone's theorem, an
exact continuous symmetry broken by the ground state produces massless
collective excitations, called \emph{Goldstone bosons}. The circular trough
illustrates one such direction; a theory can have several.\refn{8}

What plays this role in QCD? A quark field has left- and right-chiral
components. These are its two handed components; for a massless quark they
correspond to spin opposed to, or aligned with, its motion. In the idealization
of zero up- and down-quark masses, the strong-interaction equations allow
independent rotations mixing $u$ and $d$ in each handed sector. This is
\emph{chiral symmetry}.

The vacuum---the lowest-energy state of the fields---develops correlations
between these sectors. Its symmetry retains their common rotations while
breaking three independent relative directions. The associated three
collective excitations are the pions. This account describes their origin
as low-energy excitations of the same quark--gluon fields whose particle
content we introduced earlier.\refn{8}\refn{27}

The physical up and down quarks have small, nonzero masses. These make the
chiral symmetry approximate and give the pions a small restoring energy,
like a slight tilt making one point of the trough preferable. Pions
consequently have nonzero rest energies of about $135$--$140\,\mathrm{MeV}$.
In the massless-quark idealization, with electromagnetism omitted, those
pion masses vanish.\refn{26}\refn{27}

This connects back to our question about energy gaps. Adding physical light
quarks changes the low-energy spectrum; setting their masses exactly to zero
is the additional idealization that produces massless pions. For nuclear
physics, the measured small-but-positive pion mass gives the finite
$\ell_\pi\simeq1.4\,\mathrm{fm}$ range we used in the main story.

\clearpage
\section*{\textcolor{ink}{Optional deeper look: tunnelling and instantons}}

The barrier calculation in Section~2.2 used a differential equation for a
wavefunction. The path-integral formulation asks a complementary question:
which histories make the important contributions to the transition amplitude?

For a particle of mass $\mu$ moving in a smooth potential $U(x)$, the
\emph{action} of a path is the time integral of kinetic minus potential energy:
\[
S[x]=\int\left[\frac{\mu}{2}\left(\frac{\dd x}{\dd t}\right)^2-U(x)\right]\dd t.
\]
The contribution of a path carries the factor $\exp(iS[x]/\hbar)$; the
amplitude sums such contributions. Action has units of energy times time,
so its ratio to $\hbar$ is dimensionless.\refn{17}

\textbf{A change of time variable reveals an energy landscape.}
In problems where the continuation is justified, set $t=-i\tau$, with
$\tau$ real. Then $\dd t=-i\dd\tau$ and
$\dd x/\dd t=i\,\dd x/\dd\tau$. Substitution gives
\begin{align*}
S[x]
&=\int\left[-\frac{\mu}{2}\left(\frac{\dd x}{\dd\tau}\right)^2-U(x)\right](-i\,\dd\tau)\\
&=i\int\left[\frac{\mu}{2}\left(\frac{\dd x}{\dd\tau}\right)^2+U(x)\right]\dd\tau
\equiv iS_E[x].
\end{align*}
The quantity $S_E$ is the \emph{Euclidean action}. The phase factor becomes
\[
\exp(iS/\hbar)=\exp(-S_E/\hbar).
\]
Thus the continued calculation weights paths by an exponential penalty.
When the relevant action is large compared with $\hbar$, stationary paths
and nearby fluctuations provide a useful \emph{semiclassical approximation}.

For a smooth double-well potential with equal-depth minima, take $U=0$ at
the minima. A finite-action stationary path joining them in imaginary time
is an \emph{instanton}. Varying $S_E$ gives
\[
\mu\frac{\dd^2x}{\dd\tau^2}=\frac{\dd U}{\dd x}.
\]
The force has the sign of motion in the inverted potential $-U$.
An imaginary-time path can therefore connect regions separated by the
real-time barrier. Its action sets the leading exponential suppression of
the tunnelling contribution.\refn{28}

This is a useful route from elementary barrier penetration to more general
tunnelling calculations. For fusion, the path and its boundary conditions
must describe incoming nuclei at a specified collision energy and their
nuclear interactions. The double well illustrates the method; the
Coulomb-penetration problem requires its own appropriate semiclassical paths.

\textbf{An active mathematical direction: Chen's instanton density.}
For Yang--Mills fields, instantons can connect vacuum configurations
distinguished by a \emph{winding number}, an integer describing topological
wrapping. An \emph{instanton density} records local contributions to the
corresponding topological charge throughout Euclidean spacetime; integrating
it gives the total charge under suitable boundary conditions.\refn{28}\refn{29}

Jing-Yuan Chen asks how to preserve this information when the continuous
field theory is replaced by a lattice. Sampling the usual group-valued fields
loses topological distinctions. His construction refines the lattice variables
using higher category theory so that a local instanton density operator can
be defined. It provides a framework for studying correlations of this
topological structure.\refn{29} This returns us to the opening's deeper
question: how can a computable representation of quark--gluon physics retain
the mathematical structure of the underlying theory?

\clearpage
\section*{\textcolor{ink}{Open sources and the claims they support}}
\small
These notes make the boundary between derivation and imported physics explicit.
All links point to freely accessible material. The body develops its own
bookkeeping, energy cancellation, recoil estimate, and rate-counting argument;
the underlying physical laws and numerical inputs are not mathematical axioms
proved by this thesis.

\source{1}{Constituents and the meaning of quark--gluon plasma.}{US Department of
Energy, \href{https://www.energy.gov/science/doe-explainsquarks-and-gluons}{``DOE
Explains: Quarks and Gluons.''} Introductory orientation for quark content,
gluons, and the distinction between intact nucleons and deconfined matter.
The electron-cloud description does not assume classical orbital trajectories.}

\source{2}{Color charge and gluon self-interactions.}{Particle Data Group (2025),
\href{https://pdg.lbl.gov/2025/reviews/rpp2025-rev-qcd.pdf}{``Quantum
Chromodynamics,''} \S9.1, especially Eqs.~(9.1)--(9.2) and the discussion of
three- and four-gluon vertices on p.~2. These specify the microscopic theory;
we do not derive its Lagrangian or Feynman rules in the opening.}

\source{3}{Nucleon mass from dynamical QCD.}{S.~D\"urr et al.,
\href{https://arxiv.org/abs/0906.3599}{``Ab-initio Determination of Light Hadron
Masses,''} \emph{Science} \textbf{322}, 1224--1227 (2008), open manuscript.
Consult the mass-spectrum results and discussion of continuum and physical-mass
extrapolations. Quark masses and the overall physical scale are inputs; this
does not assign all mass to a single geometric flux-tube component.}

\source{4}{Asymptotic freedom.}{D.~J.~Gross,
\href{https://www.nobelprize.org/uploads/2024/10/gross-lecture-1.pdf}{``The
Discovery of Asymptotic Freedom and the Emergence of QCD,''} Nobel lecture
(2004), especially ``Implications of Asymptotic Freedom.'' For the precise
running-coupling equation, see source~2, \S9.1.1, Eq.~(9.3). A weak-coupling
result at short distance is not by itself a long-distance confinement proof.}

\source{5}{Flux tubes, a static potential, and string breaking.}{G.~S.~Bali et al.
(SESAM), \href{https://arxiv.org/abs/hep-lat/0505012}{``Observation of String
Breaking in QCD,''} \emph{Physical Review D} \textbf{71}, 114513 (2005).
The system has static quark sources and dynamical light quarks. The linear
potential is an intermediate-regime model; the final states after breaking are
color neutral. It is not a rigid-string model of a physical proton.}

\source{6}{What the unsolved mathematical problem actually asks.}{A.~Jaffe and
E.~Witten, \href{https://www.claymath.org/wp-content/uploads/2022/06/yangmills.pdf}{
``Quantum Yang--Mills Theory,''} \S4 (printed p.~6) for existence and the
spectral-gap statement; \S5 distinguishes extensions including confinement.
The full statement treats any compact simple gauge group on four-dimensional
spacetime with specified axiomatic properties. The
\href{https://www.claymath.org/millennium/yang-mills-the-maths-gap/}{Clay
problem-status page} still lists the missing proof (checked 11 September 2026).
This is not the claim that gluons acquire an ordinary particle mass.}

\begin{comment}
\source{7}{A scale-dependent-coupling analogue outside nuclei.}{P.~Coleman,
\href{https://www.physics.rutgers.edu/~coleman/articles/pwa_coleman.pdf}{
``Phil Anderson's Magnetic Ideas in Science''} (2015), \S3.5, printed p.~10.
The Kondo example supplies the established mechanism of scale-dependent
antiferromagnetic exchange and collective impurity-spin screening.
Its use as a prompt for comparing nuclear scales is the author's extrapolation.}
\end{comment}

\source{8}{The QCD-to-nuclear-force bridge.}{R.~Machleidt and D.~R.~Entem,
\href{https://arxiv.org/abs/1105.2919}{``Chiral Effective Field Theory and Nuclear
Forces,''} \emph{Physics Reports} \textbf{503}, 1--75 (2011). See the
introductory discussion and sections on the chiral expansion, pion exchange,
contact interactions, and many-nucleon forces. This is a systematic effective
description with fitted inputs and quantified truncation errors.}

\newpage
\section*{\textcolor{ink}{Open sources, continued}}
\small
\source{9}{Nuclear notation and the binding-energy trend.}{P.~Cappellaro,
\href{https://ocw.mit.edu/courses/22-02-introduction-to-applied-nuclear-physics-spring-2012/d0d046f78c917f107d925f11ac862ae4_MIT22_02S12_lec_ch1.pdf}{
MIT 22.02, Chapter 1: Introduction to Nuclear Physics} (2012), \S1.1.1 and
\S1.2. The liquid-drop explanation is qualitative here; no fitted mass formula
is used to predict light-nucleus binding energies or the exact peak location.}

\source{10}{The deliberately short weak-interaction detour.}{L.~E.~Marcucci,
R.~Schiavilla, and M.~Viviani,
\href{https://arxiv.org/abs/1303.3124}{``The Proton--Proton Weak Capture in
Chiral Effective Field Theory''} (2013). The solar reaction is not free proton
decay, nor the only physical route by which deuterium can form.}

\source{11}{Numerical binding-energy inputs.}{Atomic Mass Evaluation 2020,
\href{https://www-nds.iaea.org/amdc/ame2020/mass_1.mas20.txt}{\texttt{mass\_1.mas20}}
from the IAEA Atomic Mass Data Center. Rows $(Z,A)=(1,2),(1,3),(2,4),(26,56),
(28,62)$ provide the relevant values; the binding column is \emph{keV per
nucleon}, so multiply by $A$ and divide by $1000$ for total MeV. Atomic-mass
conventions require electronic-binding corrections for high-precision bare-nucleus
energies; they are immaterial to the $0.001\,\mathrm{MeV}$ inputs used here.}

\source{12}{Units and temperature conversion.}{NIST,
\href{https://www.nist.gov/pml/special-publication-330/sp-330-section-2}{
SI Brochure, \S2: The International System of Units}. Exact defining values:
$e=1.602176634\times10^{-19}\,\mathrm C$ and
$k_B=1.380649\times10^{-23}\,\mathrm{J\,K^{-1}}$.
Their quotient gives $k_B=8.617333262\ldots\times10^{-5}\,\mathrm{eV\,K^{-1}}$.
Together with NIST's
\href{https://physics.nist.gov/cuu/Constants/Table/allascii.txt}{CODATA table
entry for vacuum electric permittivity},
these give $e^2/(4\pi\epsilon_0)\simeq1.440\,\mathrm{MeV\,fm}$.}

\source{13}{Two-body kinematics and the meaning of cross section.}{Particle
Data Group (2025), \href{https://pdg.lbl.gov/2025/reviews/rpp2025-rev-kinematics.pdf}{
``Kinematics,''} \S49.3 (invariant amplitudes), \S49.4.2 (two-body final states),
and \S49.5 (cross sections). The main-text recoil result is a leading
nonrelativistic estimate, not an exact relativistic energy partition.}

\source{14}{The quantum law used in the barrier illustration.}{R.~P.~Feynman,
R.~B.~Leighton, and M.~Sands,
\href{https://www.feynmanlectures.caltech.edu/III_16.html}{\emph{The Feynman
Lectures on Physics}, Vol.~III, Chapter 16}, especially the Schr\"odinger-equation
discussion. Our constant-barrier ODE illustrates penetration only; it is not
used as a numerical D--T cross-section model.}

\source{15}{Amplitudes and interference.}{Feynman, Leighton, and Sands,
\href{https://www.feynmanlectures.caltech.edu/III_03.html}{Vol.~III, Chapter 3:
Probability Amplitudes}, especially the rules for combining amplitudes. The
main text uses a normalized wave-packet basis to avoid assigning a finite
probability to one exactly specified continuum momentum.}

\source{16}{A real microscopic D--T calculation and its limits.}{G.~Hupin,
S.~Quaglioni, and P.~Navr\'atil,
\href{https://www.nature.com/articles/s41467-018-08052-6}{``Ab Initio Predictions
for Polarized Deuterium--Tritium Thermonuclear Fusion,''}
\emph{Nature Communications} \textbf{10}, 351 (2019). See Results, Fig.~1,
and Methods: the ``NCSMC-pheno'' result includes a $-5\,\mathrm{keV}$ resonance
centroid correction. No polarized-fuel enhancement is assumed in our opening.}

\source{17}{Action and the sum over histories.}{Feynman, Leighton, and Sands,
\href{https://www.feynmanlectures.caltech.edu/II_19.html}{Vol.~II, Chapter 19:
The Principle of Least Action}, concluding quantum discussion. This supplies
an accessible conceptual entry, not a construction of the QCD functional measure.}

\source{18}{Later engineering connection.}{J.~P.~Freidberg, F.~J.~Mangiarotti,
and J.~Minervini,
\href{https://library.psfc.mit.edu/catalog/reports/2010/15ja/15ja021/15ja021_full.pdf}{
``Designing a Tokamak Fusion Reactor---How Does Plasma Physics Fit In?''}
(2015), MIT report PSFC/JA-15-21. This is the subsequent design target, not the
source of a confinement proof or a universal temperature threshold.}

\newpage
\section*{\textcolor{ink}{Additional sources}}
\small

\begin{comment}
\source{19}{The quantum liquid drop and induced nuclear interactions.}{
G.~Potel Aguilar and R.~A.~Broglia, \emph{The Nuclear Cooper Pair:
Structure and Reactions} (Cambridge University Press, 2022), Chapter~1,
printed pp.~1--4 for liquid-drop cohesion and surface energy, and pp.~8--17
for particle--vibration coupling and induced interactions. These supply the
physical examples used in the proposed comparison across scales.}
\end{comment}

\source{20}{BESIII's 2026 glueball result.}{BESIII Collaboration,
\href{https://arxiv.org/abs/2607.20366}{Lightest $0^{-+}$ Glueball as
Dominant Constituent of $X(2370)$} (submitted 22 July 2026).
The analysis combines mass, spin-parity, production, and decay evidence to
identify a dominant pseudoscalar-glueball component in $X(2370)$.
The result was announced at ICHEP on 5 August 2026, as reported by the
\href{https://ihep.cas.cn/xwdt2022/rd/202608/t20260805_8259263.html}{
Institute of High Energy Physics}.}

\begin{comment}
\source{21}{The historical proton--electron neutron conjecture.}{J.~Chadwick,
\href{https://www.nobelprize.org/uploads/2018/06/chadwick-lecture.pdf}{
The Neutron and Its Properties}, Nobel Lecture (12 December 1935),
opening discussion of Rutherford's 1920 conjecture.
H.~Yukawa, \href{https://www.nobelprize.org/uploads/2018/06/yukawa-lecture.pdf}{
Meson Theory in Its Developments}, Nobel Lecture (12 December 1949),
opening discussion of early attempts to explain nuclear binding through
electron or electric-charge exchange by analogy with chemical bonding.
These are the historical origins of the electromagnetic question posed here.}

\source{22}{A modern negative-cloud picture.}{J.~Friedrich and Th.~Walcher,
\href{https://arxiv.org/abs/hep-ph/0303054}{A Coherent Interpretation of the
Form Factors of the Nucleon in Terms of a Pion Cloud and Constituent Quarks},
\emph{European Physical Journal A} \textbf{17}, 607--623 (2003).
This phenomenological account describes charge structure through constituent
quarks and a pion cloud; the virtual proton--negative-pion component gives
the modern image used in the comparison.}
\end{comment}

\source{23}{Localization, momentum spread, and kinetic energy.}{Feynman,
Leighton, and Sands,
\href{https://www.feynmanlectures.caltech.edu/III_16.html}{Vol.~III,
Chapter 16: Dependence of Amplitudes on Position}, especially Eq.~(16.35)
for $\Delta x\,\Delta p_x\geq\hbar/2$, and
\href{https://www.feynmanlectures.caltech.edu/III_20.html}{Chapter 20:
Operators}. Section~0.4 combines this uncertainty relation with the variance
identity and nonrelativistic kinetic energy; the $5.2\,\mathrm{MeV}$ is a
one-coordinate localization bound for the stated width.}

\source{24}{Exponential range and the nuclear potential.}{R.~Machleidt,
\href{https://www.scholarpedia.org/article/Nuclear_Forces}{``Nuclear Forces,''}
\emph{Scholarpedia} \textbf{9}(1), 30710 (2014), subsection ``Yukawa's idea
of 1935,'' derives the exponential-over-distance form from the massive-field
equation. RIKEN's first-party account of the lattice calculation
by N.~Ishii, S.~Aoki, and T.~Hatsuda,
\href{https://www.riken.jp/en/news_pubs/research_news/rr/5066/}{
``Don't Get Too Close''} (2007), explains the attractive region, repulsive
core, and exponentially decreasing tail. The central-potential picture is
a useful description of the competition; detailed calculations include
spin and other operator structure.}

\source{25}{Identical nucleons, spin, and the exclusion principle.}{Feynman,
Leighton, and Sands,
\href{https://www.feynmanlectures.caltech.edu/III_04.html}{Vol.~III,
Chapter 4: Identical Particles}, particularly \S4--2 on the Pauli principle
and \S4--3 on nuclear scattering and spin. The orbital explanation in the
opening is an independent-particle guide to allowed occupations; the
many-body interaction sets the binding energy.}

\source{26}{Mesons and pion quark content.}{Particle Data Group (2025),
\href{https://pdg.lbl.gov/2025/reviews/rpp2025-rev-quark-model.pdf}{
``Quark Model,''} \S15.2--15.3, especially the quark-charge table and the
light-meson classification. Pion masses and the approximate
$939\,\mathrm{MeV}$ nucleon rest energy are rounded experimental scales.}

\source{27}{Why the pion mass is small.}{H.~Leutwyler,
\href{https://arxiv.org/abs/hep-ph/9609465}{``Light Quark Effective Theory''}
(1996), \S1 and \S3, especially Eq.~(3) for the leading pion-mass dependence
on the light-quark masses. Read alongside source~8, \S2.1.2--2.1.3, for
explicit and spontaneous chiral-symmetry breaking. The optional note uses
the circular-trough analogy to introduce the symmetry mechanism before its
QCD application.}

\source{28}{Imaginary-time paths and quantum-mechanical instantons.}{D.~Tong,
\href{https://davidtong.org/pdfs/teaching/gauge-theory/gauge2.pdf}{
\emph{Gauge Theory}, Chapter~2}, printed pp.~56--59 on quantum-mechanical
tunnelling in a double well. It develops the Euclidean path integral, the
instanton solution, and its semiclassical interpretation. Our optional note
shows the time-continuation algebra; prefactors require the fluctuations
about the stationary path.}

\source{29}{Preserving instanton topology on a lattice.}{J.-Y.~Chen,
\href{https://arxiv.org/abs/2406.06673v3}{``Instanton Density Operator in
Lattice QCD from Higher Category Theory,''} \emph{SciPost Physics}
\textbf{19}, 158 (2025). See \S1.1 for the topology problem and \S4.2
for the refined lattice Yang--Mills construction.}

\end{document}
